% Generated draft: requires mathematical and source review.
\documentclass[11pt]{article}
\usepackage[margin=25mm]{geometry}
\usepackage{fontspec}
\setmainfont{Latin Modern Roman}
\usepackage{amsmath,amssymb,mathtools,unicode-math}
\setmathfont{Latin Modern Math}
\usepackage{enumitem,needspace,xurl,hyperref}
\setlength{\parindent}{0pt}
\setlength{\parskip}{5pt}
\setlength{\emergencystretch}{3em}
\newcommand{\N}{\mathbb{N}}
\newcommand{\Z}{\mathbb{Z}}
\newcommand{\Q}{\mathbb{Q}}
\newcommand{\R}{\mathbb{R}}
\newcommand{\C}{\mathbb{C}}
\newcommand{\F}{\mathbb{F}}
\newcommand{\A}{\mathbb{A}}
\newcommand{\PP}{\mathbb{P}}
\newcommand{\E}{\mathbb{E}}
\newcommand{\eps}{\varepsilon}
\newcommand{\dd}{\,\mathrm{d}}
\begin{document}
\section*{UnsolvedMath ID: 30002975}
\textbf{Source problem code:} catalogue entry
\subsection{Definitions and mathematical statement}
A polarized K3 surface is a pair \((X,L)\), where \(X\) is a smooth projective K3 surface and \(L\) is an ample line bundle. Here a K3 surface has trivial canonical bundle and \(H^1(X,\mathcal O_X)=0\). The polarization is primitive if its class in \(\operatorname{Pic}(X)\) is not a positive integral multiple of another class. Its genus \(p\) is specified by \(L^2=2p-2\). Write \(|nL|\) for the complete linear system of \(L^{\otimes n}\).

Under the usual nodal-curve interpretation, let
\[
V_{n,\delta}(X,L)=\{C\in |nL|: C\text{ is reduced and irreducible, with precisely }\delta\text{ ordinary nodes and no other singularities}\}.
\]
An ordinary node means a geometric singularity whose completed local ring is isomorphic to \(\overline{k}[[x,y]]/(xy)\), after passage to an algebraic closure of the base field \(k\). The conventional numerical range is
\[
p\in\Z,\quad p\geq2,\qquad n\in\Z,\quad n\geq1,\qquad 0\leq\delta\leq p_a(nL)=1+n^2(p-1).
\]
In this interpretation the normalization of \(C\) has genus \(1+n^2(p-1)-\delta\). These numerical conventions and the exclusion of other singularities are made explicit here; the supplied record does not state them separately.

The supplied clean statement asks about the universal Severi parameter space
\[
\mathcal V_{p,n,\delta}=\{((X,L),C):(X,L)\text{ is primitively polarized of genus }p,\ C\in V_{n,\delta}(X,L)\}/\text{isomorphism},
\]
where the displayed expression describes its moduli objects, rather than specifying a scheme or stack construction. The question is whether \(\mathcal V_{p,n,\delta}\) is irreducible for every intended admissible choice of \((p,n,\delta)\) for which it is nonempty.

\emph{Unresolved qualifications.} The base field, characteristic, moduli construction, precise parameter range, and treatment of empty loci are not specified. The nonempty-locus formulation above makes an explicit convention whose agreement with the intended source question requires review. Moreover, the original statement asks whether \(V_{n,\delta}(X,L)\) is irreducible for a general polarized K3 surface \((X,L)\), whereas the clean statement asks whether the universal space \(\mathcal V_{p,n,\delta}\) is irreducible. These are distinct questions: irreducibility of the total parameter space does not by itself establish irreducibility of a general fiber. The original statement also does not explicitly impose primitivity. Both formulations are retained as unresolved alternatives; no equivalence is asserted.
\subsection{Short English statement}
The clean statement asks whether the space parametrizing a primitively polarized K3 surface together with an irreducible nodal curve in \(|nL|\) is irreducible. The original wording instead concerns curves on a general fixed polarized K3 surface. This distinction, the base field, and the exact admissible parameters require clarification. The supplied status is \emph{partially solved}; this is a dated curation label, not an independently verified assessment.
\subsection{Sources}
\begin{itemize}
\item Local UnsolvedMath registry: supplied complete source record, permanent ID \(30002975\), source problem code \texttt{catalogue entry}, in the Oberwolfach Reports Open Problems collection. No individual proposer is supplied.
\item Supplied report citation: \emph{Mini-Workshop: Singular Curves on \(K3\) Surfaces and Hyperk\"ahler Manifolds} (2016), \emph{Oberwolfach Reports}, pp.\ 2939--2967, DOI \texttt{10.4171/owr/2015/51}. \url{https://doi.org/10.4171/owr/2015/51}. The record gives proposed year 2015; the citation's year 2016 and the DOI's 2015 component are retained without reconciliation.
\item Supplied statement-verification metadata labels the statement \texttt{corrected\_verified} and reports comparison with the official Oberwolfach report. This is a supplied verification claim; no independent source consultation is asserted here.
\item Ciliberto--Dedieu, \emph{Irreducibility of Severi Varieties}. \url{https://arxiv.org/abs/1809.03914}.
\item Bruno--Lelli-Chiesa, \emph{Irreducibility of Severi Varieties}. \url{https://arxiv.org/abs/2112.09398}.
\item Ciliberto--Flamini--Galati--Knutsen, \emph{Moduli of Nodal Curves}. \url{https://arxiv.org/abs/1502.07378}.
\item Supplied literature assessment, checked on \textbf{2026-08-21}: status \emph{partially solved}. The supplied assessment reports irreducibility results in broad ranges, including \(p\geq4\delta-3\) for primitive classes and small-node ranges for general K3 surfaces, while describing the full \(\mathcal V_{p,n,\delta}\) question as open. This is supplied dated curation, not independently verified. The record expressly requests independent verification of time-sensitive status and cited claims; it does not provide the complete hypotheses or identify the precise parameter spaces covered by each cited result.
\end{itemize}
\end{document}
